\documentclass[10pt,a4paper]{amsart}
\usepackage[margin=19mm]{geometry}
\usepackage{amsmath,amssymb,mathtools}
\usepackage[scr=rsfs,cal=euler]{mathalfa}
\usepackage{xfrac}
\usepackage[unicode,hidelinks,bookmarks=false]{hyperref}
\newcommand{\ie}{i.e.\ }
\newcommand{\cf}{cf.\ }
\newcommand{\resp}{respectively\ }
\newcommand{\Subsection}{\subsection}
\providecommand{\nobreakdash}{\nobreak\hbox{-}}
\setlength{\emergencystretch}{2em}
\multlinegap0pt
\usepackage{bm}
\usepackage{bbm}
\usepackage{dsfont}
\usepackage{xspace}
\usepackage{cancel}
\usepackage{mathtools}
\usepackage{amsthm}
\makeatletter
\@ifundefined{theorem}{\newtheorem{theorem}{Theorem}}{}
\makeatother
\newcommand{\OO}{\mathcal O}
\title[Scott topologies on products of countable complete Heyting a]{Scott topologies on products of countable complete Heyting algebras}
\author{Xiaoquan Xu}
\date{Source: arXiv:2609.18032v1 (CC BY 4.0)}
\begin{document}
\maketitle

\noindent\emph{Source and licence.} Scott topologies on products of countable complete Heyting algebras. Version arXiv:2609.18032v1 (CC BY 4.0), distributed under the Creative Commons Attribution 4.0 International licence, as recorded in the arXiv licence metadata for this exact version and shown on the abstract page https://arxiv.org/abs/2609.18032v1. Full rights notice: licenses/arxiv-2609.18032-CC-BY-4.0.txt.\par\smallskip

\noindent\textit{Editorial scope.} Counted unit: the Theorem whose source label is \texttt{thm:consonance-characterization} in the article (source0.tex lines 397--416, source label \texttt{thm:consonance-characterization}), delivered together with the article's own complete original proof of it (source0.tex lines 418--477, one \texttt{proof} environment of 60 lines). No mathematics is written, repaired or re-derived: statement and proof are the article's own text line for line. Required context retained uncounted: none. Every definition, display and label of those lines is retained unchanged. Omitted from this package and not redistributed: 1--396, 417--417, 478--1575. Nothing inside a retained excerpt is omitted or altered: every retained body is delivered exactly as the article's lines are written, so no omission receipt of any kind accompanies this package. Citations: the retained excerpts carry no citation command at all, so no bibliography entry of the article is transcribed and no reference is redistributed. Reference adaptation: none was needed and none was made; every reference inside the delivered unit resolves inside it, so no unresolved reference or citation is produced. Printed slips kept literal: no printed word, symbol, hypothesis or number of the counted statement or of its proof was altered. Layout and class: the article's own class is \texttt{amsart} with packages including fontenc, lmodern, amsmath, amssymb, mathtools, microtype, needspace, natbib, hyperref; the wrapper is the lane's pinned \texttt{amsart} editorial wrapper at 10pt on A4, which restates the theorem-like environments the retained text needs and adds the article's own macro definitions that the retained text needs (\texttt{\textbackslash OO}), with extra wrapper packages (bm, bbm, dsfont, xspace, cancel, mathtools). The delivered numbering is the unit's own. Assets: no retained span contains a figure environment, a graphics-inclusion command, a PDF-page inclusion, a table or a TikZ picture; no image file is copied into this package, the package declares no asset dependency and no image file is redistributed with it. Licence: version arXiv:2609.18032v1 of this article is distributed under the Creative Commons Attribution 4.0 International licence, as recorded in the arXiv licence metadata for this exact version and shown on the abstract page https://arxiv.org/abs/2609.18032v1. A full rights notice is delivered at licenses/arxiv-2609.18032-CC-BY-4.0.txt. Characters: every retained excerpt is pure ASCII, so the delivered engine, XeLaTeX with its default Latin Modern text font, reports no missing character.
\begin{theorem}\label{thm:consonance-characterization}
Let $X_i$ $(i\in I)$ be topological spaces, and set
$X=\coprod_{i\in I}X_i$.
\begin{enumerate}
\item If $X$ is consonant, then $\Theta$ is a homeomorphism
\[
   \Sigma\OO(X)\longrightarrow\prod_{i\in I}\Sigma\OO(X_i).
\]
Equivalently,
\begin{equation}\label{eq:scott-product-open}
  \sigma\!\left(\prod_{i\in I}\OO(X_i)\right)
  =\prod_{i\in I}\sigma(\OO(X_i)).
\end{equation}
\item If every $X_i$ is consonant and \eqref{eq:scott-product-open} holds,
then $X$ is consonant.
\end{enumerate}
Consequently, for a family of consonant spaces, the Scott-product identity
\eqref{eq:scott-product-open} holds if and only if their topological sum is
consonant.
\end{theorem}

\begin{proof}
Because $\Theta$ is an order isomorphism, it identifies the Scott topology on
$\OO(X)$ with the Scott topology on the product poset
$\prod_i\OO(X_i)$.  The product topology
$\prod_i\sigma(\OO(X_i))$ is always contained in the latter Scott topology,
since every coordinate projection preserves directed suprema.

Assume first that $X$ is consonant.  Let $\mathcal H$ be Scott open in
$\OO(X)$ and $U\in\mathcal H$.  Choose a compact saturated set $K\subseteq U$
such that
\[
  U\in\Phi_X(K)\subseteq\mathcal H.
\]
Since the summands form an open cover of $X$, compactness of $K$ implies
that
\[
  F=\{i\in I:K\cap X_i\ne\varnothing\}
\]
is finite.  Put $K_i=K\cap X_i$.  Since $X_i$ is clopen in $X$, each $K_i$
is compact and saturated in $X_i$.  Then
\begin{equation}\label{eq:compact-cylinder}
  \Theta(\Phi_X(K))
  =\prod_{i\in I}\Phi_{X_i}(K_i).
\end{equation}
For $i\notin F$, $K_i=\varnothing$ and
$\Phi_{X_i}(K_i)=\OO(X_i)$; for $i\in F$, $\Phi_{X_i}(K_i)$ is Scott
open.  Hence the right-hand side of
\eqref{eq:compact-cylinder} is a basic product-open neighborhood of
$\Theta(U)$ contained in $\Theta(\mathcal H)$.  Thus $\Theta(\mathcal H)$ is
product open, proving (1).

Conversely, assume that every $X_i$ is consonant and that
\eqref{eq:scott-product-open} holds.  Let $\mathcal H$ be Scott open in
$\OO(X)$ and $U\in\mathcal H$.  Since $\Theta(\mathcal H)$ is product open,
there are a finite set $F\subseteq I$ and Scott-open families
$\mathcal H_i\subseteq\OO(X_i)$ $(i\in F)$ such that, writing
$U_i=U\cap X_i$,
\begin{equation}\label{eq:basic-neighborhood}
  \Theta(U)\in
  \left(\prod_{i\in F}\mathcal H_i\right)
  \times\left(\prod_{i\notin F}\OO(X_i)\right)
  \subseteq \Theta(\mathcal H).
\end{equation}
By consonance of $X_i$, choose compact saturated $K_i\subseteq U_i$ such
that
\[
  U_i\in\Phi_{X_i}(K_i)\subseteq\mathcal H_i
  ~ (i\in F).
\]
Let $K=\bigcup_{i\in F}K_i$, regarded as a subset of the topological sum
$X$.  Since $F$ is finite and the $X_i$ are clopen summands, $K$ is compact
and saturated in $X$; moreover, $K\subseteq U$.  If
$V\in\Phi_X(K)$, then $V\cap X_i\in\Phi_{X_i}(K_i)\subseteq\mathcal H_i$
for every $i\in F$.  Equation~\eqref{eq:basic-neighborhood} therefore gives
$V\in\mathcal H$.  Hence
\[
  U\in\Phi_X(K)\subseteq\mathcal H,
\]
and $X$ is consonant.
\end{proof}

\end{document}
